\chapter{Summary}

Peat is a common ingredient of growing media, and its removal degrades peatlands; biochar is among the growing media materials that are suggested for replacing it. The effect of biochar as a peat alternative on the salt stress response of three mint taxa, \textit{Mentha piperita}, \textit{Mentha crispata}, \textit{Mentha spicata} `Moroccan' were investigated in this study.

The experiment was conducted in a growth chamber, with 27 pots (one pot per combination), containing 3 to 4 plants with 400~g of substrate. These three treatments were: unamended peat, peat with 15~g or 30~g of the peat replaced with biochar and 0, 25 or 50~mM NaCl applied 3 times within 2 weeks. After 4 weeks from the first application, the plant height, leaf size, Fv/Fm ratio, chlorophyll content and relative water content of the leaves were measured from the upper third of the plants. Within each taxon, the effects were described and tested by two-way analysis of variance of pot means.

All taxa were reduced in height, size of leaves and chlorophyll content due to salinity. In \textit{M. piperita} plant height at 50~mM was 31\% less than at 0~mM in unamended peat; in \textit{M. crispata} plant height at 50~mM was 44\% lower than at 0~mM and in \textit{M. spicata} `Moroccan', the most sensitive taxon, it was 51\% lower at 50~mM than at 0~mM. The Fv/Fm ratio dropped just a bit. Biochar not only caused taller plants and larger leaves under all salinity levels but also mitigated the salinity-induced loss of growth and chlorophyll, especially under 30~g as compared to 15~g, and in the more sensitive taxa. The results for the relative water content were not explained and require checking.

Biochar is found to be in part a replacement of peat in the growing media of mint and to diminish the impact of salt stress. Only one pot was used for each combination so results are preliminary and should be replicated.
